\documentclass[12pt,a4paper]{article}
\usepackage[utf8]{inputenc}
\usepackage[T1]{fontenc}
\usepackage{geometry}
\geometry{margin=1in}
\usepackage{hyperref}
\usepackage{booktabs}
\usepackage{enumitem}

\title{\textbf{Project Synopsis: PassGuard AI}\\ \large Intelligent Password Strength Analyzer using Ensemble Machine Learning}
\author{Astro-Saurav}
\date{\today}

\begin{document}

\maketitle

\section{Introduction}
In the modern digital landscape, weak passwords remain the primary entry point for cyberattacks. Traditional password meters rely on basic regex patterns which are easily bypassed. \textbf{PassGuard AI} is a full-stack solution that replaces these rigid rules with an \textbf{Ensemble Machine Learning} model. It evaluates passwords across 12 distinct mathematical and linguistic features to provide a high-fidelity intelligence score.

\section{Objectives}
\begin{itemize}[noitemsep]
    \item To develop a non-deterministic password evaluation system using AI.
    \item To provide real-time, actionable security recommendations to users.
    \item To implement a secure, industry-standard authentication system using JWT.
    \item To ensure high availability and scalability through cloud-native deployment (Render/Vercel).
\end{itemize}

\section{Proposed System Architecture}
\begin{description}
    \item[Frontend:] A modern, responsive dashboard built with vanilla HTML5/CSS3/JS, utilizing asynchronous fetch API calls for a seamless user experience.
    \item[Backend:] A high-performance FastAPI server (Python) handling rate-limiting, JWT-based authentication, and ML inference.
    \item[Machine Learning Layer:] A Voting Classifier ensemble consisting of \textbf{Random Forest} and \textbf{Gradient Boosting} algorithms, trained on 15,000+ real-world and synthetic password patterns.
\end{description}

\section{Technical Specifications}
\subsection{Software Requirements}
\begin{itemize}[noitemsep]
    \item Python 3.9+ (Backend)
    \item FastAPI \& Uvicorn (Web Framework)
    \item Scikit-Learn (ML Engine)
    \item JWT (Authentication Engine)
\end{itemize}

\subsection{Hardware Requirements (Minimum)}
\begin{itemize}[noitemsep]
    \item Processor: 1.5GHz+ (Dual Core)
    \item RAM: 512MB
    \item Disk Space: 100MB
\end{itemize}

\section{Methodology}
The system operates in four primary stages:
\begin{enumerate}
    \item \textbf{Data Ingestion:} Raw password strings are captured via a secure, masked input.
    \item \textbf{Feature Extraction:} The string is decomposed into 12 features, including \textbf{Shannon Entropy}, character diversity, sequential runs, and keyboard pattern detection.
    \item \textbf{ML Inference:} The processed feature vector is passed through the pre-trained ensemble model to predict a strength class.
    \item \textbf{Actionable Feedback:} An AI-driven suggestion generator provides specific security tips based on missing entropy.
\end{enumerate}

\section{Conclusion}
PassGuard AI represents a significant shift from static security rules to dynamic, intelligent analysis. By combining modern web technologies with advanced data science, it provides a robust defense against modern password-cracking techniques.

\end{document}
